\chapter{Eigenvalues and Eigenvectors}

The determinant gives one number that says how a matrix scales volume. It does not
say what the matrix does to individual directions. This chapter answers that
question.

A matrix acts as a linear transformation: given an input vector, it produces an
output vector. Usually the output points in a new direction: the matrix may
rotate, stretch, shear, or reflect the vector. In standard coordinates, $x$ and $y$
get mixed together, which makes the geometry hard to see at a glance.

Most matrices, however, have a few special directions that they do not turn.
Along such a direction the matrix acts like multiplication by a single number: an
input vector on that line comes out as a multiple of itself (pointing the same way
or the opposite way). Those directions are the \emph{eigenvectors}, and the
multiplying number is the \emph{eigenvalue}. Once you know them, a complicated
transformation splits into separate stretches along a few lines. Engineers use
this to find how a bridge or a building vibrates, search engines use it to rank web
pages, and data scientists use it to find the main directions of spread in a data
set.

\section{Directions a matrix does not turn}
\label{ch23:directions}

Start with a picture and a single concrete matrix. Take
\[
  \mat{A}=\begin{bmatrix} 2 & 1 \\ 1 & 2 \end{bmatrix}.
\]
This is an ordinary symmetric matrix. Watch what it does to a handful of
input vectors. Apply it to $\vec{e}_1 = \colvec{1\\0}$ and you get $\colvec{2\\1}$:
the output has swung upward and points in a new direction. Apply it to
$\colvec{0\\1}$ and you get $\colvec{1\\2}$, swung the other way. For most
inputs, $\mat{A}$ turns the arrow.

But now try the diagonal direction $\vec{v}=\colvec{1\\1}$:
\[
  \mat{A}\colvec{1\\1}
  =\begin{bmatrix} 2 & 1 \\ 1 & 2 \end{bmatrix}\colvec{1\\1}
  =\colvec{2+1\\1+2}=\colvec{3\\3}=3\colvec{1\\1}.
\]
The output $\colvec{3\\3}$ points in \emph{exactly} the same
direction as the input $\colvec{1\\1}$. The matrix did not turn this arrow at all;
it tripled its length. And try the anti-diagonal $\vec{w}=\colvec{1\\-1}$:
\[
  \mat{A}\colvec{1\\-1}
  =\colvec{2-1\\1-2}=\colvec{1\\-1}=1\cdot\colvec{1\\-1}.
\]
Again the direction is preserved, this time with no stretch at all (the
length is unchanged because the multiplier is $1$).
Figure~\ref{ch23:fig-eigendirections} shows all three inputs before and after.

\begin{figure}[H]
\centering
\begin{tikzpicture}[scale=0.92,>=Latex]
  % Panel 1: Input Space
  \begin{scope}[shift={(0,0)}]
    \node[anchor=south, font=\small\bfseries, spacegray] at (1.0, 3.85) {Before: Input Space};
    \draw[help lines,gray!25] (-1.4,-1.4) grid (3.2,3.2);
    \draw[->,gray!60] (-1.4,0)--(3.4,0) node[right]{\footnotesize$x$};
    \draw[->,gray!60] (0,-1.4)--(0,3.4) node[above]{\footnotesize$y$};

    % Invariant span lines
    \draw[vecred!40,dashed,thick] (-1.2,-1.2) node[below left,vecred!80!black,font=\footnotesize] {line $L_1$} -- (3.2,3.2);
    \draw[targetgreen!40,dashed,thick] (-1.2,1.2) -- (2.2,-2.2) node[below right,targetgreen!80!black,font=\footnotesize] {line $L_2$};

    % Generic vector
    \draw[->,vecblue,very thick] (0,0) -- (1,0) node[midway,below=2pt,font=\footnotesize] {$\vec{x}$};

    % Eigenvectors
    \draw[->,vecred,very thick] (0,0) -- (1,1) node[above left=1pt,font=\footnotesize] {$\vec{v}_1$};
    \draw[->,targetgreen,very thick] (0,0) -- (1,-1) node[below left=1pt,font=\footnotesize] {$\vec{v}_2$};
  \end{scope}

  % Transformation arrow
  \draw[->,thick,gray!70] (4.0,1.2) -- (4.8,1.2) node[midway,above,font=\small] {$\mat{A}$};

  % Panel 2: Output Space
  \begin{scope}[shift={(6.0,0)}]
    \node[anchor=south, font=\small\bfseries, spacegray] at (1.0, 3.85) {After: Space Stretched};
    \draw[help lines,gray!25] (-1.4,-1.4) grid (3.4,3.4);
    \draw[->,gray!60] (-1.4,0)--(3.6,0) node[right]{\footnotesize$x$};
    \draw[->,gray!60] (0,-1.4)--(0,3.5) node[above]{\footnotesize$y$};

    % Invariant span lines
    \draw[vecred!40,dashed,thick] (-1.2,-1.2) node[below left,vecred!80!black,font=\footnotesize] {line $L_1$} -- (3.4,3.4);
    \draw[targetgreen!40,dashed,thick] (-1.2,1.2) -- (2.2,-2.2) node[below right,targetgreen!80!black,font=\footnotesize] {line $L_2$};

    % Generic vector turned off its line
    \draw[->,vecblue!60,dashed] (0,0) -- (1,0);
    \draw[->,vecblue,very thick] (0,0) -- (2,1) node[above right=1pt,font=\footnotesize] {$\mat{A}\vec{x}$};

    % Eigenvectors stretched along their lines
    \draw[->,vecred,very thick] (0,0) -- (3,3) node[above right=1pt,font=\footnotesize] {$\mat{A}\vec{v}_1=3\vec{v}_1$};
    \draw[->,targetgreen,very thick] (0,0) -- (1,-1) node[right=2pt,font=\footnotesize] {$\mat{A}\vec{v}_2=1\vec{v}_2$};
  \end{scope}
\end{tikzpicture}
\caption{The matrix $\mat{A}=\protect\begin{bmatrix}2&1\\1&2\protect\end{bmatrix}$
turns most vectors off their line (blue $\vec{x}$ turns from $(1,0)$ to $(2,1)$).
Vectors on the diagonal lines $L_1$ and $L_2$ stay on their own lines:
$\vec{v}_1$ is stretched by the factor $\lambda_1=3$, while $\vec{v}_2$ is left
unchanged ($\lambda_2=1$). The vectors on these two lines are the eigenvectors.}
\label{ch23:fig-eigendirections}
\end{figure}

\begin{intuitionbox}
An eigenvector specifies a line through the origin that is mapped into itself by the
matrix. The matrix scales vectors along that line without rotating them off it. The
eigenvalue is the scalar factor: $\lambda=3$ triples length, $\lambda=1$ leaves vectors
unchanged, $\lambda=0.5$ halves length, and $\lambda=-2$ doubles length while reversing
direction.
\end{intuitionbox}

The word \emph{eigen} is German for ``own'' or ``characteristic'' or ``proper.''
These are the matrix's own directions, which belong to the matrix itself rather
than to the coordinates we happen to use. Here is the precise statement.

\begin{definition}[Eigenvalue and eigenvector]
\label{ch23:def-eigen}
Let $\mat{A}$ be an $n\times n$ matrix. A nonzero vector $\vec{x}\in\R^n$ is an
\emph{eigenvector} of $\mat{A}$ if there is a scalar $\lambda$ with
\[
  \boxed{\;\mat{A}\vec{x}=\lambda\,\vec{x}.\;}
\]
The scalar $\lambda$ is the \emph{eigenvalue} belonging to $\vec{x}$. (We always
require $\vec{x}\neq\vec{0}$. For now, our vectors live in real space $\R^n$.
Occasionally, a real matrix forces $\lambda$ to be a complex number, as when it
rotates the plane; we explore what that means geometrically in Section~\ref{ch23:complex},
while vectors with complex entries (the space $\C^n$) wait for Chapter~32
and the Chapter~39 appendix.)
\end{definition}

Two remarks on this definition.

\textbf{Why the vector must be nonzero.} If we allowed $\vec{x}=\vec{0}$, the
equation $\mat{A}\vec{0}=\lambda\vec{0}$ would be true for \emph{every} $\lambda$,
which would make the eigenvalue meaningless. The zero vector has no direction, so
it cannot tell us a direction the matrix preserves. We throw it out.

\textbf{Why the eigenvalue \emph{can} be zero.} A legitimate eigenvalue is
$\lambda=0$. That says $\mat{A}\vec{x}=\vec{0}$ for a nonzero $\vec{x}$, i.e.
$\vec{x}$ is in the nullspace (the subspace of all vectors that $\mat{A}$ squashes to $\vec{0}$).
So ``$0$ is an eigenvalue of $\mat{A}$'' is another way of saying ``$\mat{A}$ is singular.''
Hold onto that; it will return.

\begin{focusbox}
$\mat{A}\vec{x}=\lambda\vec{x}$ is the most important equation in this half of the
book. Read it as: \emph{applying the matrix has the same effect as multiplying by
a single number}, but only along the special directions $\vec{x}$. The number is
$\lambda$; the direction is $\vec{x}$.
\end{focusbox}

\begin{remark}[An eigenvector stands for a whole line]
\label{ch23:rem-scaling}
If $\mat{A}\vec{x}=\lambda\vec{x}$, then for any scalar $c\neq 0$ the vector
$c\vec{x}$ is an eigenvector too, with the \emph{same} eigenvalue:
\[
  \mat{A}(c\vec{x})=c\,\mat{A}\vec{x}=c\lambda\vec{x}=\lambda(c\vec{x}).
\]
So eigenvectors come in lines (and, as we will see, sometimes whole planes). When we
``find the eigenvector'' we mean ``find the line''; we usually report one
convenient representative, like $\colvec{1\\1}$ rather than $\colvec{5\\5}$.
\end{remark}

\begin{recoverybox}
\textbf{Pause and think:} Suppose $\vec{v}$ is an eigenvector of a matrix $\mat{A}$ with eigenvalue $\lambda=4$. What vector does $\mat{A}$ produce when applied to $3\vec{v}$? What about $-\vec{v}$?
\vspace{0.3em}

By linearity, $\mat{A}(3\vec{v}) = 3(\mat{A}\vec{v}) = 3(4\vec{v}) = 4(3\vec{v})$, and $\mat{A}(-\vec{v}) = 4(-\vec{v})$. In both cases, $\mat{A}$ multiplies the vector by $4$. The eigenvalue belongs to the \emph{entire line} through $\vec{v}$, not only to the vector $\vec{v}$.
\end{recoverybox}

\begin{walkbox}
Before reading on, picture three transformations of the plane and ask which
directions survive untouched. (1) A pure scaling that doubles everything: which
directions are eigenvectors? (2) A reflection across the $x$-axis: which arrows
stay on their line, and what are their eigenvalues? (3) A $90^\circ$ rotation:
can \emph{any} real arrow stay on its line? Carry these on your walk; we answer
all three in this chapter.
\end{walkbox}

\section{Finding the eigenvalues: the characteristic equation}
\label{ch23:char}

Geometry told us eigenvectors exist for the matrix $\begin{bmatrix}2&1\\1&2\end{bmatrix}$
because we got lucky and guessed $\colvec{1\\1}$. We need a method that does not
rely on luck.

Start from $\mat{A}\vec{x}=\lambda\vec{x}$ and move everything to one side. The
trick is to write $\lambda\vec{x}=\lambda\mat{I}\vec{x}$ so that both terms are a
matrix times $\vec{x}$:
\begin{align*}
  \mat{A}\vec{x}-\lambda\vec{x}&=\vec{0}\\
  \mat{A}\vec{x}-\lambda\mat{I}\vec{x}&=\vec{0} &&\text{(since $\mat{I}\vec{x}=\vec{x}$)}\\
  (\mat{A}-\lambda\mat{I})\,\vec{x}&=\vec{0}. &&\text{(factor out $\vec{x}$)}
\end{align*}

\begin{trapbox}
You cannot write $\mat{A}\vec{x}-\lambda\vec{x}=(\mat{A}-\lambda)\vec{x}$ and stop
there. The symbol $\mat{A}-\lambda$ is nonsense: $\mat{A}$ is a matrix and
$\lambda$ is a number, and you cannot subtract a number from a matrix. You must
insert the identity, $\lambda\vec{x}=\lambda\mat{I}\vec{x}$, turning the scalar
into the matrix $\lambda\mat{I}=\diag(\lambda,\dots,\lambda)$. \emph{Then}
$\mat{A}-\lambda\mat{I}$ is a legitimate matrix subtraction.
\end{trapbox}

Consider $(\mat{A}-\lambda\mat{I})\vec{x}=\vec{0}$. This is a homogeneous
system, and we want a \emph{nonzero} solution $\vec{x}$. When does a square system
$\mat{M}\vec{x}=\vec{0}$ have a nonzero solution? Exactly when the matrix
$\mat{M}=\mat{A}-\lambda\mat{I}$ is \emph{singular}, that is, when it squashes some
nonzero vector down to $\vec{0}$, i.e.\ when its nullspace contains more than
$\vec{0}$. The determinant chapters gave us a test for exactly this: a square
matrix is singular precisely when its determinant is zero.

\begin{formalbox}
The number $\lambda$ is an eigenvalue of $\mat{A}$ if and only if
$\mat{A}-\lambda\mat{I}$ is singular, that is, if and only if
\[
  \boxed{\;\det(\mat{A}-\lambda\mat{I})=0.\;}
\]
The left-hand side, viewed as a function of $\lambda$, is a polynomial of degree
$n$ called the \emph{characteristic polynomial} $p(\lambda)$; the equation
$p(\lambda)=0$ is the \emph{characteristic equation}. Its roots are exactly the
eigenvalues of $\mat{A}$.
\end{formalbox}

So to find eigenvalues you do not hunt for special vectors at
all: you compute one determinant, get a polynomial in $\lambda$, and find its
roots. The eigenvectors come afterward, one eigenvalue at a time.
Figure~\ref{ch23:fig-recipe} shows the two stages.

\begin{figure}[htbp]
\centering
\begin{tikzpicture}[>=Latex,
    stage/.style={draw, rounded corners, thick, align=center, font=\small,
                  inner sep=5pt, minimum height=1.25cm},
    lab/.style={font=\small, align=center}]
  \node[stage, draw=spacegray, fill=spacegray!8] (A) at (0,0) {matrix\\ $\mat{A}$};
  \node[stage, draw=quotientorange, fill=quotientorange!10] (P) at (3.7,0)
       {$p(\lambda)=\det(\mat{A}-\lambda\mat{I})$\\ a polynomial};
  \node[stage, draw=vecred, fill=vecred!8] (L) at (7.5,0)
       {roots\\ $\lambda_1,\lambda_2,\dots$};
  \node[stage, draw=vecblue, fill=vecblue!8] (X) at (11.0,0)
       {nullspace of\\ $\mat{A}-\lambda_i\mat{I}$};
  \draw[->,thick] (A)--(P);
  \draw[->,thick] (P)--(L);
  \draw[->,thick] (L)--(X);
  \draw[decorate,decoration={brace,mirror,amplitude=6pt},thick,quotientorange!80!black]
       (0.0,-0.95)--(7.9,-0.95) node[midway,below=8pt,lab]
       {Stage 1: a determinant finds the numbers $\lambda$};
  \draw[decorate,decoration={brace,mirror,amplitude=6pt},thick,vecblue]
       (9.6,-0.95)--(12.4,-0.95) node[midway,below=8pt,lab]
       {Stage 2: elimination finds\\ the directions $\vec{x}$};
\end{tikzpicture}
\caption{The two stages of finding eigenvalues and eigenvectors. First a
determinant turns the matrix into a polynomial whose roots are the eigenvalues;
then, for each eigenvalue, solving $(\mat{A}-\lambda_i\mat{I})\vec{x}=\vec{0}$ gives
its eigenvectors.}
\label{ch23:fig-recipe}
\end{figure}

\subsection{The full recipe, walked through twice}
\label{ch23:recipe}

\begin{focusbox}
\textbf{The eigenvalue recipe.}
\begin{enumerate}
  \item Form $\mat{A}-\lambda\mat{I}$ (subtract $\lambda$ from each diagonal entry).
  \item Compute $\det(\mat{A}-\lambda\mat{I})$; this is the characteristic
        polynomial $p(\lambda)$.
  \item Solve $p(\lambda)=0$ for the eigenvalues $\lambda_1,\lambda_2,\dots$
  \item For each eigenvalue $\lambda_i$, solve $(\mat{A}-\lambda_i\mat{I})\vec{x}=\vec{0}$
        to find its eigenvectors (its nullspace).
\end{enumerate}
\end{focusbox}

\begin{example}[Our running matrix, done by the book]
\label{ch23:ex-running}
Take $\mat{A}=\begin{bmatrix}2&1\\1&2\end{bmatrix}$ and rediscover its
eigenvalues and eigenvectors without guessing.

\emph{Step 1.} Subtract $\lambda$ from the diagonal:
\[
  \mat{A}-\lambda\mat{I}=\begin{bmatrix}2-\lambda & 1\\ 1 & 2-\lambda\end{bmatrix}.
\]

\emph{Step 2.} Take the determinant of this $2\times 2$ (recall
$\det\begin{bmatrix}a&b\\c&d\end{bmatrix}=ad-bc$):
\[
  p(\lambda)=(2-\lambda)(2-\lambda)-(1)(1)=(2-\lambda)^2-1.
\]

\emph{Step 3.} Set it to zero. We could expand, but it factors as a
difference of squares:
\[
  (2-\lambda)^2-1=\bigl((2-\lambda)-1\bigr)\bigl((2-\lambda)+1\bigr)=(1-\lambda)(3-\lambda).
\]
So
\[
  p(\lambda)=(1-\lambda)(3-\lambda)=0\quad\Longrightarrow\quad
  \lambda_1=3,\ \lambda_2=1.
\]
These are the two stretch factors we found geometrically. Graphically, the
eigenvalues are where the parabola $p(\lambda)$ crosses the horizontal axis (the
left curve in Figure~\ref{ch23:fig-charpoly} on page~\pageref{ch23:fig-charpoly}).

\emph{Step 4, eigenvalue $\lambda_1=3$.} Solve
$(\mat{A}-3\mat{I})\vec{x}=\vec{0}$:
\[
  \mat{A}-3\mat{I}=\begin{bmatrix}2-3 & 1\\ 1 & 2-3\end{bmatrix}
  =\begin{bmatrix}-1 & 1\\ 1 & -1\end{bmatrix}.
\]
The single equation both rows give is $-x_1+x_2=0$, i.e. $x_1=x_2$. So every
eigenvector for $\lambda=3$ is a multiple of $\colvec{1\\1}$. (The matrix
$\begin{bmatrix}-1&1\\1&-1\end{bmatrix}$ is singular -- its two rows are
proportional -- which is exactly why $\lambda=3$ was an eigenvalue.)

\emph{Step 4, eigenvalue $\lambda_2=1$.} Solve
$(\mat{A}-1\cdot\mat{I})\vec{x}=\vec{0}$:
\[
  \mat{A}-\mat{I}=\begin{bmatrix}1 & 1\\ 1 & 1\end{bmatrix},\qquad
  x_1+x_2=0\ \Longrightarrow\ x_2=-x_1.
\]
So the eigenvectors for $\lambda=1$ are multiples of $\colvec{1\\-1}$. Both lines
match Figure~\ref{ch23:fig-eigendirections} exactly.
\end{example}

\begin{trapbox}
After you plug an eigenvalue back in, the matrix $\mat{A}-\lambda\mat{I}$ \emph{must}
come out singular (its rows must be dependent), or you have made an arithmetic
error. If you reduce $(\mat{A}-\lambda\mat{I})\vec{x}=\vec{0}$ and find the only
solution is $\vec{x}=\vec{0}$, something is wrong: either your $\lambda$ is not
really a root, or you slipped in the determinant. A correct eigenvalue always
leaves you a free variable.
\end{trapbox}

\begin{example}[A non-symmetric $2\times 2$]
\label{ch23:ex-nonsym}
Let $\mat{B}=\begin{bmatrix}4 & 1\\ 2 & 3\end{bmatrix}$. Then
\[
  \det(\mat{B}-\lambda\mat{I})
  =\det\begin{bmatrix}4-\lambda & 1\\ 2 & 3-\lambda\end{bmatrix}
  =(4-\lambda)(3-\lambda)-(1)(2).
\]
Expand: $(4-\lambda)(3-\lambda)=12-7\lambda+\lambda^2$, so
\[
  p(\lambda)=\lambda^2-7\lambda+12-2=\lambda^2-7\lambda+10=(\lambda-2)(\lambda-5).
\]
The eigenvalues are $\lambda=5$ and $\lambda=2$.

For $\lambda=5$: $\mat{B}-5\mat{I}=\begin{bmatrix}-1&1\\2&-2\end{bmatrix}$, giving
$-x_1+x_2=0$, so $\vec{x}=\colvec{1\\1}$.

For $\lambda=2$: $\mat{B}-2\mat{I}=\begin{bmatrix}2&1\\2&1\end{bmatrix}$, giving
$2x_1+x_2=0$, so $x_2=-2x_1$ and $\vec{x}=\colvec{1\\-2}$.

Check the eigenvalue $5$ by direct multiplication:
\[
  \mat{B}\colvec{1\\1}=\colvec{4+1\\2+3}=\colvec{5\\5}=5\colvec{1\\1}.
\]
\end{example}

\begin{labbox}
Compute the eigenvalues and eigenvectors of
$\mat{C}=\begin{bmatrix}3 & 0\\ 4 & 1\end{bmatrix}$ entirely by hand, following
the four steps. (Because $\mat{C}$ is triangular, you should be able to read its
eigenvalues straight off the diagonal before you even start -- see
Section~\ref{ch23:triangular}. Use the recipe anyway and confirm they agree.)
Then verify each eigenvector by direct multiplication.
\end{labbox}

\begin{strangbox}
Strang introduces eigenvalues with the same move, noting that the eigenvalues
come from $\det(\mat{A}-\lambda\mat{I})=0$ while the eigenvectors come from the
nullspace of $\mat{A}-\lambda\mat{I}$. If you are reading this \emph{before}
Strang, fix in your mind that two different ideas are stacked here: a
\emph{determinant} test finds the special numbers $\lambda$, and then a
\emph{nullspace} computation finds the special directions $\vec{x}$. Reading
\emph{after} Strang, the rest of this chapter goes past the mechanics into what the
eigenvalues \emph{mean}: trace, determinant, complex rotations.
\end{strangbox}

\section{Triangular and diagonal matrices: eigenvalues for free}
\label{ch23:triangular}

For certain structured matrices, the eigenvalues can be read directly from the entries. If $\mat{A}$ is triangular
(all zeros above, or all zeros below, the diagonal), then $\mat{A}-\lambda\mat{I}$
is also triangular, and the determinant of a triangular matrix is the product
of its diagonal entries. So
\[
  \det(\mat{A}-\lambda\mat{I})=(a_{11}-\lambda)(a_{22}-\lambda)\cdots(a_{nn}-\lambda),
\]
and the roots are $\lambda=a_{11},a_{22},\dots,a_{nn}$. (Here $a_{ii}$ is the
entry in row $i$, column $i$ of $\mat{A}$.)

\begin{proposition}
\label{ch23:prop-triangular}
The eigenvalues of a triangular (in particular, a diagonal) matrix are exactly its
diagonal entries.
\end{proposition}

\begin{example}
\label{ch23:ex-tri}
For $\mat{T}=\begin{bmatrix}5 & 8 & -2\\ 0 & 1 & 7\\ 0 & 0 & 3\end{bmatrix}$ the
eigenvalues are $5,1,3$ at a glance -- no determinant expansion needed. (The
eigenvectors still require work; only the eigenvalues are free.)
\end{example}

\begin{trapbox}
``Eigenvalues are the diagonal entries'' is true \emph{only} for triangular
matrices. For a full matrix the diagonal entries are not the eigenvalues. For
Example~\ref{ch23:ex-nonsym}, $\mat{B}$ had diagonal $4,3$ but eigenvalues $5,2$.
What \emph{is} always true is that the diagonal entries \emph{sum} to the same
thing the eigenvalues sum to -- the trace -- which is the subject of the next
section ($4+3=7=5+2$).
\end{trapbox}

\begin{recoverybox}
\textbf{Pause and think:} A $3\times 3$ diagonal matrix has entries $7, -2, 4$ along its main diagonal. Without writing down any characteristic polynomial, what are its eigenvalues, and what are its eigenvectors?
\vspace{0.3em}

Because a diagonal matrix is triangular, Proposition~\ref{ch23:prop-triangular} tells us the eigenvalues are the diagonal entries: $\lambda_1=7$, $\lambda_2=-2$, and $\lambda_3=4$. Call the matrix $\mat{D}$. Its eigenvectors are the standard unit basis vectors
\[
  \vec{e}_1 = \colvec{1\\0\\0},\qquad \vec{e}_2 = \colvec{0\\1\\0},\qquad \vec{e}_3 = \colvec{0\\0\\1},
\]
because $\mat{D}\vec{e}_1 = 7\vec{e}_1$, $\mat{D}\vec{e}_2 = -2\vec{e}_2$, and $\mat{D}\vec{e}_3 = 4\vec{e}_3$. A diagonal matrix is already written in its own eigenvector axes.
\end{recoverybox}

\section{Trace and determinant: the eigenvalues' fingerprints}
\label{ch23:trace-det}

Here are two facts that look like coincidences in our examples and turn out to be
laws. In Example~\ref{ch23:ex-nonsym} the eigenvalues were $5$ and $2$. Their
\emph{sum} is $7$, which is $4+3$, the sum of the diagonal entries of $\mat{B}$.
That diagonal sum is called the \emph{trace} of the matrix:
$\tr\mat{B} = 4+3 = 7$. In general, for any $n\times n$ matrix $\mat{A}$, the trace is
defined as the sum of all entries on its main diagonal:
\[
  \tr\mat{A} = a_{11} + a_{22} + \dots + a_{nn}.
\]
Their \emph{product} is $10$, which is exactly $\det\mat{B}=4\cdot 3-1\cdot 2$.
Both identities hold for every square matrix.

\begin{theorem}[Trace and determinant from eigenvalues]
\label{ch23:thm-trace-det}
Let $\mat{A}$ be $n\times n$ with eigenvalues $\lambda_1,\dots,\lambda_n$ listed
with algebraic multiplicity (meaning each eigenvalue is repeated as many times
as it appears as a root of $p(\lambda)$, counting complex roots). Then
\[
  \lambda_1+\lambda_2+\cdots+\lambda_n=\tr\mat{A},
  \qquad
  \lambda_1\lambda_2\cdots\lambda_n=\det\mat{A}.
\]
The eigenvalues \emph{add up} to the trace and \emph{multiply} to the determinant.
\end{theorem}

Here is why, in the $2\times 2$ case. Write
$\mat{A}=\begin{bmatrix}a&b\\c&d\end{bmatrix}$. Its characteristic polynomial is
\begin{align*}
  \det(\mat{A}-\lambda\mat{I})
  &=(a-\lambda)(d-\lambda)-bc\\
  &=\lambda^2-(a+d)\lambda+(ad-bc)\\
  &=\lambda^2-(\tr\mat{A})\,\lambda+\det\mat{A}.
\end{align*}
So for any $2\times 2$ matrix the characteristic polynomial is
\[
  \boxed{\;p(\lambda)=\lambda^2-(\tr\mat{A})\,\lambda+\det\mat{A}.\;}
\]
Now compare with the factored form. If the roots are $\lambda_1,\lambda_2$, then
\[
  p(\lambda)=(\lambda-\lambda_1)(\lambda-\lambda_2)
  =\lambda^2-(\lambda_1+\lambda_2)\lambda+\lambda_1\lambda_2.
\]
Matching the coefficient of $\lambda$ gives
$\lambda_1+\lambda_2=\tr\mat{A}$, and matching the constant term gives
$\lambda_1\lambda_2=\det\mat{A}$.

The general $n\times n$ proof uses the same idea. The characteristic polynomial
factors as
\[
  p(\lambda)=(-1)^n(\lambda-\lambda_1)(\lambda-\lambda_2)\cdots(\lambda-\lambda_n).
\]
Expanding the product, the coefficient of $\lambda^{n-1}$ is
$(-1)^{n-1}(\lambda_1+\cdots+\lambda_n)$; expanding the determinant instead, the same
coefficient comes out as $(-1)^{n-1}\tr\mat{A}$. The constant term is
$p(0)=\det(\mat{A}-0\,\mat{I})=\det\mat{A}$, and from the product it is
$\lambda_1\lambda_2\cdots\lambda_n$.

\begin{focusbox}
For a $2\times 2$ you can often skip the algebra entirely. You only need two
numbers: $t=\tr\mat{A}$ and $d=\det\mat{A}$. The eigenvalues are the two numbers
that add to $t$ and multiply to $d$:
\[
  \lambda=\frac{t\pm\sqrt{t^2-4d}}{2}.
\]
\end{focusbox}

The number under the square root, $t^2-4d$, is called the \emph{discriminant}. Its
sign tells you what kind of eigenvalues to expect before you compute them.
Figure~\ref{ch23:fig-charpoly} graphs $p(\lambda)$ for three matrices from this
chapter. If $t^2-4d>0$ the parabola crosses the axis twice: two different real
eigenvalues. If $t^2-4d=0$ it touches the axis once: one repeated eigenvalue. If
$t^2-4d<0$ it never reaches the axis: no real eigenvalues, and
Section~\ref{ch23:complex} explains what the two roots become.

\begin{figure}[htbp]
\centering
\begin{tikzpicture}
  \pgfplotsset{charpoly/.style={width=4.6cm, height=4.6cm,
      axis lines=middle, xmin=-0.6, xmax=4.2, ymin=-1.6, ymax=4.2,
      xtick={1,2,3}, ytick=\empty, samples=80, domain=-0.5:4.1,
      xlabel={$\lambda$}, ylabel={$p(\lambda)$},
      every axis x label/.style={at={(current axis.right of origin)},anchor=west},
      every axis y label/.style={at={(axis cs:0,4.2)},anchor=south east},
      tick label style={font=\small}, title style={font=\small, align=center, yshift=10pt}}}
  \begin{axis}[charpoly, name=a1,
      title={$\mat{A}$: $(\lambda-1)(\lambda-3)$\\[2pt] two real roots}]
    \addplot[vecred, very thick] {(x-1)*(x-3)};
    \addplot[only marks, mark=*, vecred] coordinates {(1,0) (3,0)};
  \end{axis}
  \begin{axis}[charpoly, name=a2, at={(a1.east)}, anchor=west, xshift=1.1cm,
      title={shear $\mat{S}$: $(\lambda-1)^2$\\[2pt] one repeated root}]
    \addplot[quotientorange, very thick] {(x-1)^2};
    \addplot[only marks, mark=*, quotientorange] coordinates {(1,0)};
  \end{axis}
  \begin{axis}[charpoly, name=a3, at={(a2.east)}, anchor=west, xshift=1.1cm,
      title={rotation $\mat{Q}$: $\lambda^2+1$\\[2pt] no real roots}]
    \addplot[structurepurple, very thick] {x^2+1};
  \end{axis}
\end{tikzpicture}
\caption{The characteristic polynomial $p(\lambda)=\lambda^2-t\lambda+d$ of three
$2\times 2$ matrices from this chapter: the running matrix
$\mat{A}=\protect\begin{bmatrix}2&1\\1&2\protect\end{bmatrix}$, the shear $\mat{S}$
and the $90^\circ$ rotation $\mat{Q}$ (both met later in this chapter). Real
eigenvalues are the points where the
graph meets the horizontal axis (dots). Left: two crossings, eigenvalues $1$ and
$3$. Middle: the graph only touches the axis, so the eigenvalue $1$ is repeated.
Right: the graph stays above the axis, so the eigenvalues are not real numbers.}
\label{ch23:fig-charpoly}
\end{figure}

\begin{example}[Reading eigenvalues off trace and determinant]
\label{ch23:ex-tracedet}
For $\mat{A}=\begin{bmatrix}2&1\\1&2\end{bmatrix}$: $t=4$, $d=4-1=3$. We need two
numbers adding to $4$ and multiplying to $3$. By inspection: $3$ and $1$. (Or use
the formula: $\lambda=\tfrac{4\pm\sqrt{16-12}}{2}=\tfrac{4\pm 2}{2}=3,1$.) Same
answer as Example~\ref{ch23:ex-running}, with no determinant of $\mat{A}-\lambda\mat{I}$
expanded.
\end{example}

\begin{aestheticbox}
The trace and determinant are computed directly from the entries, with no
eigenvalues in sight, yet they are completely determined by the eigenvalues. The
next chapter shows that the eigenvalues do not change when you describe the same
transformation in a different coordinate system (a different basis), so trace and
determinant do not change either. They act as two \emph{fingerprints} of the
transformation. The eigenvalues are the full fingerprint; the trace (their sum)
and determinant (their product) are the two you can read off without solving the
whole polynomial. That is why they make quick checks on an eigenvalue computation.
\end{aestheticbox}

\begin{labbox}
Use only trace and determinant -- do not expand a characteristic polynomial -- to
find the eigenvalues of each:
\[
  \begin{bmatrix}1&2\\2&1\end{bmatrix},\qquad
  \begin{bmatrix}0&-1\\1&0\end{bmatrix},\qquad
  \begin{bmatrix}3&0\\0&3\end{bmatrix},\qquad
  \begin{bmatrix}5&4\\0&5\end{bmatrix}.
\]
For each, the eigenvalues add to the trace and multiply to the determinant. Watch
what happens to the second one (trace $0$, determinant $1$): the discriminant
$t^2-4d=-4$ is negative. We deal with that in Section~\ref{ch23:complex}.
\end{labbox}

\section{What the eigenvalues do to the picture}
\label{ch23:geometry}

Each eigenvalue also has a direct geometric meaning. Once you know the eigenvalues
of a $2\times 2$ matrix and its eigenvector directions, you can picture what it
does to the whole plane. The method is to measure vectors along the eigenvectors
instead of along the $x$- and $y$-axes.

Suppose $\mat{A}$ has eigenvectors $\vec{v}_1,\vec{v}_2$ that span the plane, with
eigenvalues $\lambda_1,\lambda_2$. Any vector $\vec{x}$ can be written as a
combination $\vec{x}=c_1\vec{v}_1+c_2\vec{v}_2$. Then by linearity
\[
  \mat{A}\vec{x}
  =c_1\,\mat{A}\vec{v}_1+c_2\,\mat{A}\vec{v}_2
  =c_1\lambda_1\vec{v}_1+c_2\lambda_2\vec{v}_2.
\]
The matrix acts \emph{one eigendirection at a time}: it stretches the $\vec{v}_1$
part by $\lambda_1$ and the $\vec{v}_2$ part by $\lambda_2$, independently.
Figure~\ref{ch23:fig-decompose} shows this for our running matrix. Measured in
eigenvector coordinates, the matrix does nothing more than two separate scalings.
That is the starting point for \emph{diagonalization} in the next chapter:
rewriting a matrix, in eigenvector coordinates, as a diagonal matrix.

\begin{figure}[htbp]
\centering
\begin{tikzpicture}[scale=0.85,>=Latex]
  % LEFT: x = c1 v1 + c2 v2 with v1=(1,1), v2=(1,-1), c1=0.75, c2=0.5
  \begin{scope}
    \node[font=\small\bfseries, spacegray] at (1.2,3.6) {Before: $\vec{x}=c_1\vec{v}_1+c_2\vec{v}_2$};
    \draw[help lines,gray!25] (-1.2,-1.6) grid (3.6,3.2);
    \draw[->,gray!60] (-1.2,0)--(3.8,0);
    \draw[->,gray!60] (0,-1.6)--(0,3.3);
    \draw[vecred!35,dashed] (-1,-1)--(3.1,3.1);
    \draw[targetgreen!45,dashed] (-1,1)--(1.5,-1.5);
    % c1 v1 = (0.75,0.75), c2 v2 = (0.5,-0.5), x = (1.25,0.25)
    \draw[->,vecred,very thick] (0,0)--(0.75,0.75) node[above left=1pt,font=\small]{$c_1\vec{v}_1$};
    \draw[->,targetgreen,very thick] (0,0)--(0.5,-0.5) node[below left=1pt,font=\small]{$c_2\vec{v}_2$};
    \draw[gray,dashed] (0.75,0.75)--(1.25,0.25)--(0.5,-0.5);
    \draw[->,vecblue,very thick] (0,0)--(1.25,0.25) node[right=2pt,font=\small]{$\vec{x}$};
  \end{scope}
  \draw[->,thick,gray!70] (4.1,1.0)--(5.0,1.0) node[midway,above,font=\small]{$\mat{A}$};
  % RIGHT: A x = 3 c1 v1 + 1 c2 v2 = (2.25,2.25)+(0.5,-0.5) = (2.75,1.75)
  \begin{scope}[xshift=6.4cm]
    \node[font=\small\bfseries, spacegray] at (1.2,3.6) {After: $\mat{A}\vec{x}=3c_1\vec{v}_1+1c_2\vec{v}_2$};
    \draw[help lines,gray!25] (-1.2,-1.6) grid (3.6,3.2);
    \draw[->,gray!60] (-1.2,0)--(3.8,0);
    \draw[->,gray!60] (0,-1.6)--(0,3.3);
    \draw[vecred!35,dashed] (-1,-1)--(3.1,3.1);
    \draw[targetgreen!45,dashed] (-1,1)--(1.5,-1.5);
    \draw[->,vecred,very thick] (0,0)--(2.25,2.25) node[above left=1pt,font=\small]{$3c_1\vec{v}_1$};
    \draw[->,targetgreen,very thick] (0,0)--(0.5,-0.5) node[below left=1pt,font=\small]{$c_2\vec{v}_2$};
    \draw[gray,dashed] (2.25,2.25)--(2.75,1.75)--(0.5,-0.5);
    \draw[->,vecblue,very thick] (0,0)--(2.75,1.75) node[right=2pt,font=\small]{$\mat{A}\vec{x}$};
  \end{scope}
\end{tikzpicture}
\caption{Splitting a vector along the eigenvectors of
$\mat{A}=\protect\begin{bmatrix}2&1\\1&2\protect\end{bmatrix}$ (here $\vec{v}_1=(1,1)$
with $\lambda_1=3$, $\vec{v}_2=(1,-1)$ with $\lambda_2=1$). Left: $\vec{x}$ is the sum
of a red piece on the first eigenvector line and a green piece on the second.
Right: $\mat{A}$ triples the red piece, leaves the green piece alone, and adds them
back up to give $\mat{A}\vec{x}$.}
\label{ch23:fig-decompose}
\end{figure}

Figure~\ref{ch23:fig-ellipse} shows two more cases: two positive eigenvalues turn
the unit circle into a stretched ellipse, while a negative eigenvalue flips one
direction end-to-end.

\begin{figure}[H]
\centering
\begin{tikzpicture}[scale=0.95,>=Latex]
  % LEFT: lambda = 3, 1 (stretch along (1,1), keep (1,-1))
  \begin{scope}
    \draw[help lines,gray!25] (-2.4,-2.4) grid (2.4,2.4);
    \draw[->] (-2.5,0)--(2.6,0); \draw[->] (0,-2.5)--(0,2.6);
    % unit circle of inputs, image is ellipse along (1,1)
    \draw[vecblue,thick] (0,0) circle (1);
    \draw[targetgreen,thick,rotate=45] (0,0) ellipse (3 and 1);
    % eigendirections: stretched along (1,1) by 3, fixed along (1,-1) by 1
    \draw[->,vecred,very thick] (0,0)--({3/sqrt(2)},{3/sqrt(2)}) node[above left=2pt,font=\small] {$\lambda_1=3$};
    \draw[->,structurepurple,very thick] (0,0)--({1/sqrt(2)},{-1/sqrt(2)}) node[below right=3pt,font=\small] {$\lambda_2=1$};
    \node at (0,-2.95) {\small inputs (circle) $\mapsto$ stretched (ellipse)};
    \node at (0,2.95) {\small $\lambda_1=3,\ \lambda_2=1$};
  \end{scope}
  % RIGHT: lambda = 2, -1 (flip + stretch)
  \begin{scope}[xshift=7.2cm]
    \draw[help lines,gray!25] (-1.6,-1.6) grid (2.2,1.6);
    \draw[->] (-1.7,0)--(2.4,0); \draw[->] (0,-1.7)--(0,1.8);
    \draw[->,vecred,very thick] (0,0)--(1,0) node[above=2pt]{$\vec{v}_1$};
    \draw[->,structurepurple,very thick,dashed] (0,0)--(0,1) node[above left=2pt]{$\vec{v}_2$};
    \draw[->,vecred,thick] (0,0)--(2,0) node[below right=1pt]{$\lambda_1\vec{v}_1=2\vec{v}_1$};
    \draw[->,structurepurple,thick] (0,0)--(0,-1) node[below left=2pt]{$\lambda_2\vec{v}_2=-\vec{v}_2$};
    \node at (0.5,2.95) {\small $\lambda_1=2,\ \lambda_2=-1$ (one flip)};
    \node at (0.5,-2.95) {\small positive stretches, negative flips};
  \end{scope}
\end{tikzpicture}
\caption{Left: when both eigenvalues are positive, the matrix turns the unit
circle into an ellipse whose axes are the eigenvector directions, stretched by the
eigenvalues. Right: a negative eigenvalue ($\lambda_2=-1$) flips that
eigendirection end-to-end while the positive one ($\lambda_1=2$) stretches.}
\label{ch23:fig-ellipse}
\end{figure}

Here is what a single real eigenvalue $\lambda$ does to vectors on its line
(Figure~\ref{ch23:fig-lambda-line} summarizes the cases):

\begin{itemize}
  \item $\lambda>1$: \emph{stretch outward} along that line (the direction is
        amplified). Repeatedly applying $\mat{A}$ sends those components flying
        off to infinity.
  \item $\lambda=1$: \emph{fixed direction} -- vectors on that line do not move at
        all.
  \item $0<\lambda<1$: \emph{shrink} toward the origin along that line. Repeated
        application collapses these components to nothing.
  \item $\lambda=0$: \emph{collapse} -- the entire line is sent to the origin in
        one step. The matrix is singular; this line is the nullspace.
  \item $\lambda<0$: \emph{flip} across the origin, and stretch (if
        $\lambda<-1$) or shrink (if $-1<\lambda<0$). Each application of $\mat{A}$
        bounces that component to the other side.
\end{itemize}

\begin{figure}[htbp]
\centering
\begin{tikzpicture}[>=Latex, x=2.6cm]
  % shaded regions on the lambda axis
  \fill[trapred!12]        (-2.4,0) rectangle (-1,0.5);
  \fill[quotientorange!14] (-1,0)   rectangle (0,0.5);
  \fill[targetgreen!14]    (0,0)    rectangle (1,0.5);
  \fill[vecblue!12]        (1,0)    rectangle (2.4,0.5);
  \draw[->,thick] (-2.45,0)--(2.5,0) node[right]{$\lambda$};
  \foreach \t in {-1,0,1}{
    \draw[thick] (\t,0.08)--(\t,-0.08) node[below,font=\small]{$\t$};
  }
  \node[font=\small, align=center] at (-1.7,0.25) {flip and stretch};
  \node[font=\small, align=center] at (-0.5,0.25) {flip and shrink};
  \node[font=\small, align=center] at (0.5,0.25)  {shrink};
  \node[font=\small, align=center] at (1.7,0.25)  {stretch};
  % markers for special values
  \node[font=\small, align=center, below=14pt] at (0,0) {collapse\\ to $\vec{0}$};
  \node[font=\small, align=center, below=14pt] at (1,0) {fixed};
  \node[font=\small, align=center, below=14pt] at (-1,0) {pure flip};
  % example iterates for each region: vector lengths after repeated application
  \node[font=\small, align=center, below=14pt] at (-1.75,0) {e.g.\ $\lambda=-2$:\\ $1,-2,4,-8,\dots$};
  \node[font=\small, align=center, below=14pt] at (1.75,0) {e.g.\ $\lambda=2$:\\ $1,2,4,8,\dots$};
  \draw[decorate,decoration={brace,amplitude=6pt},thick,spacegray]
       (-0.98,0.6)--(0.98,0.6) node[midway,above=7pt,font=\small]
       {$|\lambda|<1$: dies away when repeated};
  \draw[decorate,decoration={brace,amplitude=6pt},thick,spacegray]
       (1.02,0.6)--(2.4,0.6) node[midway,above=7pt,font=\small]
       {$|\lambda|>1$: grows};
  \draw[decorate,decoration={brace,amplitude=6pt},thick,spacegray]
       (-2.4,0.6)--(-1.02,0.6) node[midway,above=7pt,font=\small]
       {$|\lambda|>1$: grows};
\end{tikzpicture}
\caption{What a real eigenvalue $\lambda$ does to vectors on its eigenvector line,
by where $\lambda$ sits on the number line. Negative values flip the vector to the
other side of the origin; values with $|\lambda|<1$ shrink it, values with
$|\lambda|>1$ stretch it. The numbers under the axis show the multiplier on a
vector after $0,1,2,3$ applications of the matrix.}
\label{ch23:fig-lambda-line}
\end{figure}

\begin{intuitionbox}
The size $|\lambda|$ decides what happens to each eigendirection in the long run,
when the matrix is applied over and over. The eigenvalue with the largest
$|\lambda|$ wins: after many applications of $\mat{A}$, most vectors end up
pointing almost along the eigenvector of the largest $|\lambda|$, because that
component grows fastest (or shrinks slowest) compared with the others. This fact
is behind the \emph{power method} (find the dominant eigenvector by multiplying a
starting vector by $\mat{A}$ again and again), behind PageRank below, and behind
the study of systems that evolve step by step by a rule
$\vec{x}_{k+1}=\mat{A}\vec{x}_k$. We make it precise in Chapters~24 and~35.
\end{intuitionbox}

\begin{applicationbox}
\textbf{PageRank: the web's dominant eigenvector.} PageRank is the method Google
used to rank web pages. It builds a giant matrix $\mat{M}$ describing a
``random surfer'': start on a page, follow an outgoing link chosen at random, repeat. The
entry $M_{ij}$ is the probability of stepping from page $j$ to page $i$. The
importance of every page is the eigenvector of $\mat{M}$ belonging to the
eigenvalue $\lambda=1$: the one direction that the surfing process leaves
\emph{unchanged}, $\mat{M}\vec{x}=\vec{x}$. Its entries give the long-run fraction
of time the surfer spends on each page (the \emph{steady state}). Every other
eigendirection has $|\lambda|<1$ and dies away as the surfer keeps walking, so no
matter where you start, you converge to that eigenvector. PageRank comes
down to one instruction: ``find the eigenvector for $\lambda=1$ of a matrix with a few
billion rows.'' Two facts from this chapter make it work: eigenvalue $1$ means a
direction the matrix fixes, and the largest $|\lambda|$ dominates the long run.
\end{applicationbox}

\begin{example}[Reading a matrix from its eigendata]
\label{ch23:ex-readoff}
Suppose someone tells you a $2\times 2$ matrix has eigenvalue $2$ along
$\colvec{1\\0}$ and eigenvalue $-1$ along $\colvec{0\\1}$. You can sketch its
action immediately: it doubles everything horizontal and flips everything vertical.
The matrix must be $\begin{bmatrix}2&0\\0&-1\end{bmatrix}$, because that diagonal
matrix has exactly those eigenvalues along the axes. Knowing $n$ independent
eigenvectors of an $n\times n$ matrix and their eigenvalues pins the matrix down
completely.
\end{example}

\section{Special matrices and their eigenvalues}
\label{ch23:special}

A few geometric transformations you already know have eigenvalues you can read off
from the geometry, with no algebra at all.

\textbf{Projections.} Let $\mat{P}$ project the plane onto a line $L$ through the
origin. If $\vec{x}$ already lies on $L$, projecting does nothing:
$\mat{P}\vec{x}=\vec{x}$, eigenvalue $1$. If $\vec{x}$ is perpendicular to $L$,
projecting kills it: $\mat{P}\vec{x}=\vec{0}=0\cdot\vec{x}$, eigenvalue $0$. So
every projection has eigenvalues $1$ and $0$ -- which is exactly why projections
satisfy $\mat{P}^2=\mat{P}$ (the eigenvalues satisfy $\lambda^2=\lambda$).
Figure~\ref{ch23:fig-projection} shows the two eigendirections.

\begin{figure}[H]
\centering
\begin{tikzpicture}[scale=1.0,>=Latex]
  \draw[gray!55,thick] (-1.7,-1.7)--(1.7,1.7) node[right]{\small line $L$};
  \draw[->] (-1.8,0)--(1.8,0); \draw[->] (0,-1.8)--(0,1.8);
  % a generic vector x = (1.5, 0.3) projects onto L (y=x): proj = 0.9(1,1)
  \coordinate (X) at (1.5,0.3);
  \coordinate (PX) at (0.9,0.9);
  \draw[->,vecblue,very thick] (0,0)--(X) node[right=2pt]{$\vec{x}$};
  \draw[->,vecred,very thick] (0,0)--(PX) node[above left]{$\mat{P}\vec{x}$};
  \draw[gray!60,dashed] (X)--(PX);
  % on L: lambda=1
  \draw[->,targetgreen,thick] (0,0)--(-1.1,-1.1) node[left=3pt]{on $L$: $\lambda=1$};
  % perpendicular: lambda=0 (collapses to origin)
  \draw[->,structurepurple,thick] (0,0)--(-0.85,0.85) node[above left]{$\perp L$: $\lambda=0$};
  \filldraw[structurepurple] (0,0) circle (1.6pt);
\end{tikzpicture}
\caption{Projection onto the line $L$. Vectors already on $L$ are left alone
(eigenvalue $1$); vectors perpendicular to $L$ are collapsed to the origin
(eigenvalue $0$). A generic $\vec{x}$ (blue) is dropped onto its shadow
$\mat{P}\vec{x}$ on $L$ (red), so it is moved off its own line and is not an eigenvector.}
\label{ch23:fig-projection}
\end{figure}

\textbf{Reflections.} Let $\mat{R}$ reflect across a line $L$. Vectors on $L$ are
untouched (eigenvalue $1$); vectors perpendicular to $L$ are sent to their
negatives (eigenvalue $-1$). A reflection always has eigenvalues $+1$ and $-1$,
which is why doing it twice is the identity: $\lambda^2=1$.
Figure~\ref{ch23:fig-reflection} shows the two eigendirections.

\begin{figure}[H]
\centering
\begin{tikzpicture}[scale=1.0,>=Latex]
  \begin{scope}
    \draw[gray!50,thick] (-1.6,-1.6)--(1.6,1.6) node[right]{\small line $L$};
    \draw[->] (-1.7,0)--(1.7,0); \draw[->] (0,-1.7)--(0,1.7);
    % reflect across y=x: (a,b) -> (b,a)
    \draw[->,vecblue,very thick] (0,0)--(1.4,0.4) node[right]{$\vec{x}$};
    \draw[->,vecred,very thick] (0,0)--(0.4,1.4) node[above]{$\mat{R}\vec{x}$};
    \draw[gray!60,dashed] (1.4,0.4)--(0.4,1.4);
    \draw[->,targetgreen,thick] (0,0)--(-1.0,-1.0) node[left=3pt]{on $L$: $\lambda=1$};
    \draw[->,structurepurple!80!black,thick] (0,0)--(0.7,-0.7);
    \draw[->,structurepurple!80!black,thick,dashed] (0,0)--(-0.7,0.7) node[above left]{$\perp L$: $\lambda=-1$};
    \node at (0,-2.1){\small reflection across $L$};
  \end{scope}
\end{tikzpicture}
\caption{A reflection across the line $L$ fixes the direction along $L$ (eigenvalue
$+1$) and flips the perpendicular direction (eigenvalue $-1$). A generic vector
$\vec{x}$ (blue) is turned to $\mat{R}\vec{x}$ (red), off its own line, so it is not an
eigenvector.}
\label{ch23:fig-reflection}
\end{figure}

\textbf{Scalings.} The matrix $c\mat{I}$ stretches every direction by $c$. So
\emph{every} nonzero vector is an eigenvector, all with the same eigenvalue $c$.
This is the one case where eigenvectors do not pick out special lines: the
matrix treats every direction the same way.

\textbf{Shears.} The shear $\mat{S}=\begin{bmatrix}1&1\\0&1\end{bmatrix}$ slides
the plane horizontally. Its only eigenvalue is $1$ (it is triangular, both diagonal
entries $1$), and the only eigenvectors lie along the horizontal axis
$\colvec{1\\0}$, the layer that does not slide. A shear is the classic example
of a matrix with a \emph{repeated} eigenvalue but \emph{not enough} eigenvectors
to fill the plane (only one line, not two; see Figure~\ref{ch23:fig-shear}). A
square matrix whose eigenvectors cannot be combined into a basis of the whole space
is called \emph{defective}. We will care a great deal about that
shortfall when we ask which matrices can be diagonalized.

\begin{figure}[H]
\centering
\begin{tikzpicture}[scale=1.15,>=Latex]
  \draw[help lines,gray!25] (-0.5,-0.5) grid (2.5,1.7);
  \draw[->,gray!60] (-0.5,0)--(2.7,0) node[right]{\footnotesize$x$};
  \draw[->,gray!60] (0,-0.5)--(0,1.8) node[above]{\footnotesize$y$};

  % Original unit square
  \draw[vecblue!40,dashed,thick] (0,0) -- (1,0) -- (1,1) -- (0,1) -- cycle;

  % Sheared parallelogram
  \fill[quotientorange!12] (0,0) -- (1,0) -- (2,1) -- (1,1) -- cycle;
  \draw[quotientorange,thick] (0,0) -- (1,0) -- (2,1) -- (1,1) -- cycle;

  % The only eigen-line: y = 0
  \draw[targetgreen!60,dashed,very thick] (-0.4,0) -- (2.5,0);
  \draw[->,targetgreen,very thick] (0,0) -- (1,0) node[below=3pt,font=\footnotesize] {$\mat{S}\vec{e}_1 = \vec{e}_1$};

  % Tilting vectors: e_2 and its image
  \draw[->,vecblue!60,dashed,thick] (0,0) -- (0,1) node[above left=1pt,font=\footnotesize] {$\vec{e}_2$};
  \draw[->,vecblue,very thick] (0,0) -- (1,1) node[above=2pt,font=\footnotesize] {$\mat{S}\vec{e}_2=(1,1)$};

  % Diagonal vector (1,1) and its image (2,1)
  \draw[->,vecred,thick] (0,0) -- (2,1) node[below right=1pt,font=\footnotesize] {$\mat{S}(1,1)=(2,1)$};
\end{tikzpicture}
\caption{The shear matrix $\mat{S}=\protect\begin{bmatrix}1&1\\0&1\protect\end{bmatrix}$
slides horizontal lines to the right by an amount proportional to height.
The horizontal axis (green) is the \emph{only} line preserved: $\mat{S}\colvec{1\\0}=1\cdot\colvec{1\\0}$,
so $\lambda=1$ is the sole eigenvalue. Every other line tilts (for instance, vertical $\vec{e}_2$ swings
to $(1,1)$, and diagonal $(1,1)$ swings to $(2,1)$).
Because no second independent line of eigenvectors exists in the plane, the shear is \emph{defective}.}
\label{ch23:fig-shear}
\end{figure}

\begin{trapbox}
A repeated eigenvalue does \emph{not} guarantee a repeated supply of eigenvectors.
The shear has $\lambda=1$ twice but only a one-dimensional line of eigenvectors.
Contrast $\mat{I}$, which also has $\lambda=1$ twice but a whole plane of them.
``How many times $\lambda$ is a root'' (its \emph{algebraic multiplicity}) and
``the dimension of its space of eigenvectors'' (its \emph{geometric multiplicity})
are different numbers; the second can be smaller. This gap is the reason later
chapters study defective matrices and the \emph{Jordan form}, the closest a
defective matrix can get to diagonal. (It is named after Camille Jordan and has nothing
to do with the Gauss--Jordan elimination of Chapter~8, named after Wilhelm Jordan.)
\end{trapbox}

\section{Complex eigenvalues: when a matrix is really a rotation}
\label{ch23:complex}

Return to the walking question about a $90^\circ$ rotation. Can a rotation have a
real eigenvector? Picture it: a rotation turns \emph{every} arrow by a quarter
turn. There is no direction in the plane that stays on its own line. So a $90^\circ$
rotation has \emph{no real eigenvectors at all}. Yet the characteristic equation
still hands us two roots, as soon as we allow complex numbers (numbers built
from $i=\sqrt{-1}$, described below): every degree-$2$ polynomial has two complex
roots. Those roots are complex, and they are telling us ``rotation,''
not ``stretch.''

Before computing, here is a short geometric picture of complex numbers. (We
develop vectors with complex entries in Chapter~32 and the Chapter~39 appendix;
here you need only this picture.)

\begin{intuitionbox}
\textbf{A geometric picture of complex numbers.}
On the ordinary number line, multiplying by $-1$ reverses direction: it is a
$180^\circ$ half-turn. What operation, performed \emph{twice in a row}, produces a
$180^\circ$ half-turn? A $90^\circ$ quarter-turn. That is the geometric meaning of the
imaginary unit $i$, the number with $i^2=-1$: \emph{multiplying by $i$ rotates the
plane counterclockwise by $90^\circ$}.

A \emph{complex number} has the form $z = a + bi$ with $a,b$ real; $a$ is its
\emph{real part} and $b$ its \emph{imaginary part}. Plot $z$ as the point $(a,b)$
in a plane, the \emph{complex plane}: real parts along the horizontal axis
(labelled Re), imaginary parts along the vertical axis (labelled Im).
Its distance from the origin is its \emph{magnitude} $r = |z| = \sqrt{a^2+b^2}$, and
its direction is an angle $\theta$ measured counterclockwise from the positive real
axis. So $a=r\cos\theta$ and $b=r\sin\theta$. The standard shorthand for the point
on the unit circle at angle $\theta$ is
\[
  e^{i\theta} = \cos\theta + i\sin\theta
\]
(Euler's formula), so every complex number can be written $z = r e^{i\theta}$.
Multiplying by $z = r e^{i\theta}$ does two things at once:
\begin{enumerate}
  \item It \emph{stretches} lengths by the factor $r = |z|$.
  \item It \emph{rotates} by the angle $\theta$.
\end{enumerate}
When a matrix rotates the plane instead of only stretching it along fixed lines, its
eigenvalues cannot be real numbers. They are complex numbers whose
magnitude $r$ records the stretch factor and whose angle $\theta$ records the
angle of rotation.
\end{intuitionbox}

Here is the computation. The $90^\circ$ counterclockwise rotation is
\[
  \mat{Q}=\begin{bmatrix}0 & -1\\ 1 & 0\end{bmatrix}.
\]
Its trace is $0$ and its determinant is $0\cdot 0-(-1)(1)=1$, so
\[
  p(\lambda)=\lambda^2-(\tr\mat{Q})\lambda+\det\mat{Q}=\lambda^2+1=0
  \quad\Longrightarrow\quad \lambda=\pm i.
\]
The eigenvalues are the imaginary numbers $+i$ and $-i$. Their magnitude is
$|{\pm i}|=1$: the rotation does not change any length, and the magnitude of a
complex eigenvalue records exactly that scaling. The \emph{angle} of the eigenvalue
records the rotation: $i=e^{i\pi/2}$ sits at angle $90^\circ$, which is the angle
$\mat{Q}$ rotates by. The two numbers $i$ and $-i$ are mirror images across the real
axis; such a pair $a+bi$ and $a-bi$ is called a \emph{conjugate pair}, and
$\bar z = a-bi$ is the \emph{complex conjugate} of $z=a+bi$.
Figure~\ref{ch23:fig-rotation} shows both pictures at once:
the turned arrows on the left, the conjugate pair $\pm i$ on the unit circle on the
right.

\begin{figure}[H]
\centering
\begin{tikzpicture}[scale=1.1,>=Latex]
  \begin{scope}
    \draw[help lines,gray!25] (-1.3,-1.3) grid (1.3,1.3);
    \draw[->] (-1.4,0)--(1.5,0) node[right]{$x$};
    \draw[->] (0,-1.4)--(0,1.5) node[above]{$y$};
    \foreach \a in {0,45,...,315}{
      \draw[->,vecblue,thick] (0,0)--(\a:1);
      \draw[->,vecred!70,thin] (\a:1)--({\a+90}:1);
    }
    \node[align=center] at (0,-2.05){\small every arrow turns $90^\circ$;\\ \small none stays on its line};
  \end{scope}
  \begin{scope}[xshift=5.6cm]
    \draw[->] (-1.4,0)--(1.5,0) node[right]{Re};
    \draw[->] (0,-1.4)--(0,1.5) node[above]{Im};
    \draw[gray!40] (0,0) circle (1);
    \filldraw[vecred] (0,1) circle (1.6pt) node[above right]{$\lambda=i$};
    \filldraw[vecred] (0,-1) circle (1.6pt) node[below right]{$\lambda=-i$};
    \draw[->,structurepurple,thick] (0,0)--(0,1);
    \coordinate (Aw) at (1,0);
    \coordinate (Ow) at (0,0);
    \coordinate (Bw) at (0,1);
    \draw pic[draw,angle radius=10mm,"$90^\circ$"] {angle=Aw--Ow--Bw};
    \node[align=center] at (0,-2.05){\small eigenvalues lie on\\ \small the unit circle};
  \end{scope}
\end{tikzpicture}
\caption{A $90^\circ$ rotation has no real eigendirection (left): every arrow is
turned. Its eigenvalues are complex, $\pm i$, sitting on the unit circle of the
complex plane (right; horizontal axis = real part, vertical axis = imaginary part).
The magnitude $|\lambda|=1$ says ``no stretching''; the
angle of $\lambda$ says ``rotate by $90^\circ$.''}
\label{ch23:fig-rotation}
\end{figure}

The general rotation by angle $\theta$,
\[
  \mat{Q}_\theta=\begin{bmatrix}\cos\theta & -\sin\theta\\ \sin\theta & \cos\theta\end{bmatrix},
  \qquad
  \tr=2\cos\theta,\quad \det=1,
\]
has characteristic polynomial $\lambda^2-2\cos\theta\,\lambda+1=0$, whose roots are
\[
  \lambda=\cos\theta\pm\sqrt{\cos^2\theta-1}=\cos\theta\pm i\sin\theta
  =e^{\pm i\theta}.
\]
The eigenvalues are $e^{\pm i\theta}$: magnitude $1$ (no stretch, $\det=1$), and
their angle is exactly the rotation angle $\theta$.

\begin{focusbox}
\textbf{The complex eigenvalue dictionary.} A pair of complex eigenvalues
$\lambda=r e^{\pm i\theta}$ means the matrix \emph{spirals}: it rotates by the
angle $\theta$ while scaling by the radius $r=|\lambda|$ per application.
\begin{itemize}
  \item $|\lambda|=1$: pure rotation, lengths preserved.
  \item $|\lambda|<1$: inward spiral; repeated application pulls every vector in
        to the origin.
  \item $|\lambda|>1$: outward spiral; repeated application pushes every nonzero
        vector away from the origin.
\end{itemize}
Real eigenvalues stretch along fixed lines. A conjugate pair of complex eigenvalues
rotates and scales within a plane that the matrix maps to itself (for a
$2\times 2$ matrix, the whole plane).
\end{focusbox}

\begin{example}[A spiral matrix]
\label{ch23:ex-spiral}
Consider $\mat{A}=\begin{bmatrix}1 & -1\\ 1 & 1\end{bmatrix}$. Trace $2$,
determinant $1\cdot 1-(-1)(1)=2$, so
$p(\lambda)=\lambda^2-2\lambda+2=0$ and
\[
  \lambda=\frac{2\pm\sqrt{4-8}}{2}=1\pm i.
\]
Write $1\pm i=\sqrt{2}\,e^{\pm i\pi/4}$: magnitude $\sqrt 2$, angle $45^\circ$. So
$\mat{A}$ rotates by $45^\circ$ and stretches by $\sqrt2$ each time you apply it:
an outward spiral. Indeed $\mat{A}=\sqrt2\,\mat{Q}_{45^\circ}$, which you can check
directly since
\[
  \mat{Q}_{45^\circ}=\frac{1}{\sqrt2}\begin{bmatrix}1&-1\\1&1\end{bmatrix}.
\]
Repeatedly applying $\mat{A}$ to any nonzero vector traces a spiral winding out to
infinity, $45^\circ$ at a time (Figure~\ref{ch23:fig-spiral}).
\end{example}

\begin{figure}[H]
\centering
\begin{tikzpicture}[scale=0.62,>=Latex]
  \draw[help lines,gray!20] (-5,-3) grid (3,5);
  \draw[->] (-5.2,0)--(3.3,0) node[right]{$x$};
  \draw[->] (0,-3.2)--(0,5.3) node[above]{$y$};
  % orbit of (1,0) under A = sqrt2 * Q45: radius x sqrt2, angle +45 each step
  % n: point
  \coordinate (p0) at (1,0);
  \coordinate (p1) at (1,1);
  \coordinate (p2) at (0,2);
  \coordinate (p3) at (-2,2);
  \coordinate (p4) at (-4,0);
  \draw[vecblue,thick,-{Latex}] (p0)--(p1);
  \draw[vecblue,thick,-{Latex}] (p1)--(p2);
  \draw[vecblue,thick,-{Latex}] (p2)--(p3);
  \draw[vecblue,thick,-{Latex}] (p3)--(p4);
  \foreach \p/\n in {p0/0,p1/1,p2/2,p3/3,p4/4}{
    \filldraw[vecred] (\p) circle (2.4pt);
  }
  \node[below right] at (p0) {\small $\vec{x}$};
  \node[above right=1pt] at (p1) {\small $\mat{A}\vec{x}$};
  \node[above right] at (p2) {\small $\mat{A}^2\vec{x}$};
  \node[above left] at (p3) {\small $\mat{A}^3\vec{x}$};
  \node[below left] at (p4) {\small $\mat{A}^4\vec{x}$};
\end{tikzpicture}
\caption{The points $\vec{x},\mat{A}\vec{x},\mat{A}^2\vec{x},\dots$ for
$\vec{x}=\protect\colvec{1\\0}$ and
the spiral matrix $\mat{A}=\protect\begin{bmatrix}1&-1\\1&1\protect\end{bmatrix}$.
Each step turns the vector $45^\circ$ counterclockwise (the angle of the complex
eigenvalue) and lengthens it by $\sqrt2$ (the magnitude $|\lambda|>1$), so the
points wind outward. Complex eigenvalues with $|\lambda|>1$ always produce an
outward spiral like this.}
\label{ch23:fig-spiral}
\end{figure}

\begin{trapbox}
When the discriminant $t^2-4d$ is negative, do not declare ``no eigenvalues.''
There are still exactly two eigenvalues; they are complex conjugates. For a
\emph{real} matrix, complex eigenvalues always come in conjugate pairs
$\lambda,\bar\lambda$, because the characteristic polynomial has real coefficients
(the entries are real), and the complex roots of a polynomial with real
coefficients come in conjugate pairs. So a real
$2\times 2$ has either two real eigenvalues or one conjugate pair, never exactly
one complex eigenvalue alone.
\end{trapbox}

\begin{connectionbox}
\textbf{Eigenvalues, differential equations, and stability.} Eigenvalues also
describe systems that change continuously in time. Consider the differential
equation $\frac{d\vec{x}}{dt}=\mat{A}\vec{x}$, which says the rate of change of
$\vec{x}(t)$ at each moment is $\mat{A}$ times its current value. For each
eigenvalue $\lambda$ with eigenvector $\vec{v}$ (an \emph{eigenpair}
$(\lambda,\vec{v})$), the function $e^{\lambda t}\vec{v}$ is a solution, because
\[
  \frac{d}{dt}\bigl(e^{\lambda t}\vec{v}\bigr)=\lambda e^{\lambda t}\vec{v}
  =e^{\lambda t}\mat{A}\vec{v}=\mat{A}\bigl(e^{\lambda t}\vec{v}\bigr).
\]
The real part of $\lambda$, written $\operatorname{Re}\lambda$, decides the long-run
size: $\operatorname{Re}\lambda<0$ means $e^{\lambda t}\to 0$, a \emph{stable} mode
that decays; $\operatorname{Re}\lambda>0$ means exponential growth. The imaginary
part sets how fast the solution oscillates. A bridge
oscillates, a circuit rings, a population cycles, and each can be read off from where the
eigenvalues sit in the complex plane. (We develop this fully in Chapter~25.
The same complex eigenvalues that meant ``rotation'' for repeated multiplication
by $\mat{A}$ mean ``oscillation'' for this differential equation.)
\end{connectionbox}

\section{Why eigenvectors of distinct eigenvalues are independent}
\label{ch23:independence}

We keep saying ``the eigenvectors span the plane'' as if it were automatic. It is
not automatic, as the shear matrix showed, but distinct eigenvalues always guarantee it.

\begin{theorem}[Distinct eigenvalues give independent eigenvectors]
\label{ch23:thm-indep}
If $\vec{v}_1,\dots,\vec{v}_k$ are eigenvectors of $\mat{A}$ with \emph{distinct}
eigenvalues $\lambda_1,\dots,\lambda_k$ (all different from one another), then
$\vec{v}_1,\dots,\vec{v}_k$ are linearly independent.
\end{theorem}

\begin{proof}
Take the smallest case $k=2$ first; it shows the idea. Suppose
$c_1\vec{v}_1+c_2\vec{v}_2=\vec{0}$. Apply $\mat{A}$ to both sides; since
$\mat{A}\vec{v}_i=\lambda_i\vec{v}_i$,
\[
  c_1\lambda_1\vec{v}_1+c_2\lambda_2\vec{v}_2=\vec{0}.
\]
Now also multiply the original relation by $\lambda_1$:
$c_1\lambda_1\vec{v}_1+c_2\lambda_1\vec{v}_2=\vec{0}$. Subtract the two:
\[
  c_2(\lambda_2-\lambda_1)\vec{v}_2=\vec{0}.
\]
Because $\vec{v}_2\neq\vec{0}$ and $\lambda_2\neq\lambda_1$, the only way out is
$c_2=0$. Back-substituting, $c_1\vec{v}_1=\vec{0}$ forces $c_1=0$ as well. So the
only combination giving $\vec{0}$ is the trivial one: $\vec{v}_1,\vec{v}_2$ are
independent. The general case is the same trick by induction: assume any $k-1$ of
them are independent; a dependence among all $k$, hit with $\mat{A}$ and combined
with a $\lambda_k$-multiple of itself, collapses to a dependence among $k-1$, which
forces all the coefficients to vanish.
\end{proof}

\begin{corollary}
\label{ch23:cor-basis}
If an $n\times n$ matrix has $n$ \emph{distinct} eigenvalues, its $n$ eigenvectors
form a basis of $\R^n$ (of $\C^n$ if some eigenvalues are complex). Every vector then decomposes into
eigendirections, and the matrix acts as $n$ independent scalings.
\end{corollary}

A basis made of eigenvectors of $\mat{A}$ is called an \emph{eigenbasis}. So
$n$ distinct eigenvalues always produce an eigenbasis. The shear matrix of
Section~\ref{ch23:special} has no eigenbasis, and the cause is its repeated
eigenvalue (though, as the $3\times 3$ example below shows, a repeated eigenvalue
does not always cause this). Which matrices have an eigenbasis is the central
question of diagonalization in the next chapter.

\begin{strangbox}
Strang states the central conclusion of this chapter directly: eigenvectors from
different eigenvalues are independent, so a matrix with no repeated eigenvalues has
a full set of eigenvectors. If you are heading \emph{into} Strang, carry
Theorem~\ref{ch23:thm-indep} as the reason diagonalization works. If you are
coming \emph{out} of Strang, this is the theorem you will lean on to write
$\mat{A}=\mat{S}\mat{\Lambda}\mat{S}\inv$ in the next chapter (the columns of
$\mat{S}$ are the eigenvectors, and $\mat{\Lambda}$ is the diagonal matrix of
eigenvalues), and to justify the spectral theorem for symmetric matrices later on
(the result that a symmetric matrix has real eigenvalues and perpendicular
eigenvectors).
\end{strangbox}

\section{A worked $3\times 3$}
\label{ch23:threebythree}

Nothing about the recipe changes in three dimensions; the determinant is a
little bigger. Here is one worked all the way through.

\begin{example}[Eigenstructure of a $3\times 3$]
\label{ch23:ex-3x3}
Take
\[
  \mat{A}=\begin{bmatrix}2 & 0 & 0\\ 1 & 2 & 1\\ -1 & 0 & 1\end{bmatrix}.
\]
\emph{Characteristic polynomial.} Expand $\det(\mat{A}-\lambda\mat{I})$ along the
first row, which has two zeros and so is easy:
\[
  \det\begin{bmatrix}2-\lambda & 0 & 0\\ 1 & 2-\lambda & 1\\ -1 & 0 & 1-\lambda\end{bmatrix}
  =(2-\lambda)\det\begin{bmatrix}2-\lambda & 1\\ 0 & 1-\lambda\end{bmatrix}.
\]
The remaining $2\times 2$ is triangular, so its determinant is
$(2-\lambda)(1-\lambda)$. Altogether
\[
  p(\lambda)=(2-\lambda)\,(2-\lambda)(1-\lambda)=(2-\lambda)^2(1-\lambda).
\]
The eigenvalues are $\lambda=2$ (a double root) and $\lambda=1$.

\emph{Check with trace and determinant.} The trace is $2+2+1=5$, and the
eigenvalues sum to $2+2+1=5$. Good. The determinant is
$\lambda_1\lambda_2\lambda_3=2\cdot2\cdot1=4$; you can verify directly that
$\det\mat{A}=4$ by cofactor expansion. Both fingerprints match.

\emph{Eigenvectors for $\lambda=1$.} Solve $(\mat{A}-\mat{I})\vec{x}=\vec{0}$:
\[
  \mat{A}-\mat{I}=\begin{bmatrix}1 & 0 & 0\\ 1 & 1 & 1\\ -1 & 0 & 0\end{bmatrix}.
\]
Row one says $x_1=0$; row two then says $x_2+x_3=0$, so $x_2=-x_3$. Taking
$x_3=1$ gives the eigenvector $\colvec{0\\-1\\1}$.

\emph{Eigenvectors for $\lambda=2$.} Solve $(\mat{A}-2\mat{I})\vec{x}=\vec{0}$:
\[
  \mat{A}-2\mat{I}=\begin{bmatrix}0 & 0 & 0\\ 1 & 0 & 1\\ -1 & 0 & -1\end{bmatrix}.
\]
The first row is all zeros; the second and third rows both reduce to
$x_1+x_3=0$. So there are \emph{two} free variables ($x_2$ free, and $x_3$ free
with $x_1=-x_3$). The eigenvectors for $\lambda=2$ fill a whole \emph{plane},
spanned by $\colvec{0\\1\\0}$ (taking $x_2=1,x_3=0$) and $\colvec{-1\\0\\1}$
(taking $x_2=0,x_3=1$). The subspace formed by all eigenvectors belonging to an
eigenvalue $\lambda$, together with the zero vector $\vec{0}$, is called the
\emph{eigenspace} $E_\lambda = \operatorname{null}(\mat{A}-\lambda\mat{I})$:
it is the set of all directions that $\mat{A}$ scales by $\lambda$.

Although $\lambda=2$ is a root of multiplicity two, its eigenspace is two-dimensional.
Combining this two-dimensional eigenspace with the one-dimensional eigenspace for $\lambda=1$
gives three linearly independent eigenvectors spanning all of $\R^3$. This matrix has an
eigenbasis,
\[
  \left\{\colvec{0\\-1\\1},\ \colvec{0\\1\\0},\ \colvec{-1\\0\\1}\right\},
\]
despite the repeated eigenvalue. In the shear matrix a repeated eigenvalue produced
only a one-dimensional eigenspace; here the geometric multiplicity ($2$) equals the
algebraic multiplicity ($2$).
Figure~\ref{ch23:fig-3x3} pictures the result: a single eigen-line for $\lambda=1$ and an entire
eigen-plane for $\lambda=2$, together filling $\R^3$.
\end{example}

\begin{figure}[H]
\centering
\begin{tikzpicture}[scale=1.0,>=Latex,line join=round]
  % the lambda=2 eigen-plane spanned by (0,1,0) and (-1,0,1):
  \fill[targetgreen!18,draw=targetgreen!60]
    (-1.05,1.85)--(1.35,1.35)--(0.05,-1.85)--(-2.35,-1.35)--cycle;

  % simple oblique 3D axes
  \def\ax{(2.6,0)}     % x to the right
  \def\ay{(-1.1,-0.75)}% y receding
  \def\az{(0,2.6)}     % z up
  \draw[->,gray!70] (0,0)--\ax node[right]{$x_1$};
  \draw[->,gray!70] (0,0)--\ay node[left=3pt]{$x_2$};
  \draw[->,gray!70] (0,0)--\az node[above]{$x_3$};

  % spanning eigenvectors inside the plane
  \draw[->,vecblue,very thick] (0,0)--(-1.1,-0.75) node[below right=2pt,font=\footnotesize]{\small $(0,1,0)$};
  \draw[->,vecblue,very thick] (0,0)--(-1.3,1.3)
        node[above left=1pt,font=\footnotesize]{\small $(-1,0,1)$};
  % the lambda=1 eigen-line through (0,-1,1): screen dir = -ay + az
  \draw[vecred!50,dashed,thick] (-0.8,-0.55-1.9) -- (0,0);
  \draw[->,vecred,very thick] (0,0)--(0.9,0.61+2.1) node[above right,font=\small] {$\lambda=1$ line};
  \node[above right,text=targetgreen!80!black] at (1.35,1.35) {\small $\lambda=2$ plane};
  \filldraw (0,0) circle (1.5pt);
\end{tikzpicture}
\caption{Eigenspaces of the $3\times 3$ matrix of Example~\ref{ch23:ex-3x3}. The
double eigenvalue $\lambda=2$ has an entire plane of eigenvectors (green), spanned by
$(0,1,0)$ and $(-1,0,1)$ (blue), while $\lambda=1$ has a single
eigen-line (red) along $(0,-1,1)$. The plane and the line meet only at the origin
and together span $\R^3$, so the matrix has an eigenbasis.}
\label{ch23:fig-3x3}
\end{figure}


\begin{labbox}
By hand, find all eigenvalues of
$\mat{A}=\begin{bmatrix}2 & 1 & 0\\ 0 & 2 & 0\\ 0 & 0 & 3\end{bmatrix}$
(triangular -- read them off), then determine the dimension of the eigenspace for
$\lambda=2$ by row-reducing $\mat{A}-2\mat{I}$. You should find $\lambda=2$ is a
double root but its eigenspace is only one-dimensional: a $3\times 3$ cousin of the
shear, a matrix that \emph{cannot} be diagonalized. Compare with
Example~\ref{ch23:ex-3x3}, where the double root was diagonalizable.
\end{labbox}

\begin{historybox}
Eigenvalues and eigenvectors did not arise from abstract algebra; they came
from physical problems in eighteenth-century astronomy and mechanics. When
Leonhard Euler (1750s) studied the rotation of spinning solid bodies, he
found that describing the motion in arbitrary coordinates produced a tangle of
coupled equations. He proved that every rigid body has
three mutually perpendicular \emph{principal axes of inertia}: the
directions about which an object can spin steadily without wobbling.
Three decades later, Joseph-Louis Lagrange (1780s) tackled celestial stability:
do the mutual gravitational attractions of the planets eventually destabilize the
solar system? In standard coordinates, the acceleration of every planet
depends on every other planet. Lagrange sought a new coordinate system that would
split this coupled system into independent \emph{harmonic
oscillators} (simple back-and-forth periodic motions, like ideal pendulums).

In 1829, Augustin-Louis Cauchy recast these ideas in terms of what we now call
matrices, studying the equation $\det(\mat{A}-\lambda\mat{I})=0$ (later named the
characteristic equation) and proving that real symmetric matrices always have
real eigenvalues. Decades
later, David Hilbert popularized the German term \emph{Eigenwert} (meaning ``own
value'' or ``characteristic value''). Two hurdles caught early mathematicians.
The first was the defective matrix: early geometers assumed every linear transformation
had a full set of invariant axes, and shear matrices, with only a single line of
eigenvectors, showed otherwise. The second was computational: because
eigenvalues are polynomial roots, mathematicians hoped for explicit algebraic formulas,
but by the Abel--Ruffini theorem (1824), polynomial equations of degree five or higher
cannot in general be solved by radicals (formulas using only $+,-,\times,\div$ and
root extractions like $\sqrt[n]{\cdot}$). So no finite formula gives the eigenvalues
of a general $5\times 5$ or larger matrix; they must be computed by an iterative
process that gets closer with each step. Around 1960, Soviet mathematician Vera
Kublanovskaya and British computer scientist John Francis independently invented
the QR algorithm, which repeatedly applies the QR factorization of
Chapter~20 to drive a matrix toward triangular form, where
the eigenvalues can be read off the diagonal.

Eigenvalues identify the directions of a transformation along which
coupled dynamics split into one-dimensional scalings, and they determine
long-term growth, decay, or steady-state behavior. Their main
limitations are that they apply only to square matrices, and that the eigenvectors
can fail to span the space. In Chapter~24, we assemble these directions
into a basis to diagonalize matrices, in Chapter~25 we use them to solve
differential equations, and in Chapter~29 we meet the singular value decomposition, which extends
the idea of stretching axes to rectangular matrices.
\end{historybox}

\section{Pulling it together}
\label{ch23:together}

\begin{aestheticbox}
The relation $\mat{A}\vec{x}=\lambda\vec{x}$ picks out the directions
along which a matrix acts by a single number. The eigenvalues summarize the
transformation in a way that does not depend on coordinates: their sum is the trace,
their product is the determinant, real eigenvalues describe stretching along lines,
and complex eigenvalues describe rotation. Their sizes decide what happens when the
matrix is applied over and over. Next comes \emph{diagonalization}
($\mat{A}=\mat{S}\mat{\Lambda}\mat{S}\inv$), where the eigenvectors form a coordinate system
in which the transformation is diagonal, and later the spectral theorem and the
singular value decomposition.
\end{aestheticbox}

\begin{focusbox}
\textbf{What to hold onto from this chapter.}
\begin{itemize}
  \item $\mat{A}\vec{x}=\lambda\vec{x}$ with $\vec{x}\neq\vec{0}$: eigenvectors are
        the directions a matrix does not turn; eigenvalues are how much it scales
        them.
  \item Eigenvalues are the roots of $\det(\mat{A}-\lambda\mat{I})=0$; eigenvectors
        are the nullspace of $\mat{A}-\lambda\mat{I}$.
  \item For $2\times 2$: eigenvalues add to the trace and multiply to the
        determinant; $p(\lambda)=\lambda^2-(\tr\mat{A})\lambda+\det\mat{A}$.
  \item Triangular matrices show their eigenvalues on the diagonal.
  \item $\lambda=0\Leftrightarrow$ singular; $|\lambda|$ governs long-run growth;
        complex eigenvalues $r e^{\pm i\theta}$ mean rotate-by-$\theta$,
        scale-by-$r$.
  \item Distinct eigenvalues $\Rightarrow$ independent eigenvectors $\Rightarrow$
        a complete eigenbasis (the doorway to diagonalization).
  \end{itemize}
\end{focusbox}

\section*{Exercises}
\addcontentsline{toc}{section}{Exercises}

\begin{exercise}[Quick mental checks]
\label{ch23:ex-quick-checks}
Answer each with a brief reason, without setting up any equations:
\begin{enumerate}[label=\textbf{(\alph*)}]
  \item Can the zero vector $\vec{0}$ ever be an eigenvector? Can the number $0$ ever be an eigenvalue?
  \item If a $2\times 2$ matrix has eigenvalues $\lambda_1 = 4$ and $\lambda_2 = -3$, what are its trace and its determinant?
  \item If $\mat{A}\vec{v} = 5\vec{v}$, what is $\mat{A}^3\vec{v}$? What is $(\mat{A}+2\mat{I})\vec{v}$?
  \item Can a $3\times 3$ matrix with all real entries have three purely imaginary eigenvalues? (\emph{Hint:} complex roots of real polynomials come in conjugate pairs.)
\end{enumerate}
\end{exercise}

\begin{exercise}[Warmup by hand]
\label{ch23:ex-warm}
For each matrix, find the eigenvalues and one eigenvector per eigenvalue.
\[
  \text{(a) } \begin{bmatrix}3 & 0\\ 0 & -2\end{bmatrix}\quad
  \text{(b) } \begin{bmatrix}2 & 1\\ 0 & 2\end{bmatrix}\quad
  \text{(c) } \begin{bmatrix}0 & 2\\ 2 & 0\end{bmatrix}\quad
  \text{(d) } \begin{bmatrix}1 & 4\\ 1 & 1\end{bmatrix}.
\]
\end{exercise}

\begin{exercise}[Trace and determinant only]
\label{ch23:ex-tracedet-only}
Without expanding any characteristic polynomial, find the eigenvalues of
$\begin{bmatrix}6 & 2\\ 2 & 3\end{bmatrix}$ using only that they add to the trace
and multiply to the determinant. Then verify by computing
$\det(\mat{A}-\lambda\mat{I})$ directly.
\end{exercise}

\begin{exercise}[Reflection geometry]
\label{ch23:ex-reflect}
The matrix $\mat{R}=\begin{bmatrix}0 & 1\\ 1 & 0\end{bmatrix}$ reflects the plane
across the line $y=x$. \textbf{(a)} Without computing, predict its eigenvalues and
eigenvectors from the geometry. \textbf{(b)} Confirm with the characteristic
equation. \textbf{(c)} Explain why $\mat{R}^2=\mat{I}$ follows from the eigenvalues.
\end{exercise}

\begin{exercise}[A scaled rotation]
\label{ch23:ex-rot}
Let $\mat{A}=\begin{bmatrix}3 & -2\\ 2 & 3\end{bmatrix}$. \textbf{(a)} Show the
eigenvalues are complex and write them in the form $r e^{\pm i\theta}$. \textbf{(b)}
By what angle does $\mat{A}$ rotate, and by what factor does it scale, each time it
is applied? \textbf{(c)} What is the long-run fate of $\mat{A}^n\vec{x}$ for a
generic $\vec{x}$ -- does it spiral in or out?
\end{exercise}

\begin{exercise}[Singular means $\lambda=0$]
\label{ch23:ex-singular}
Show that the matrix $\begin{bmatrix}2 & 4\\ 1 & 2\end{bmatrix}$ has $0$ as an
eigenvalue, find the corresponding eigenvector, and explain how this eigenvector
relates to the nullspace and to the fact that $\det\mat{A}=0$.
\end{exercise}

\begin{exercise}[Eigenvalues of $\mat{A}^2$ and $\mat{A}\inv$]
\label{ch23:ex-powers}
Suppose $\mat{A}\vec{x}=\lambda\vec{x}$ with $\vec{x}\neq\vec{0}$. \textbf{(a)} Show
$\mat{A}^2\vec{x}=\lambda^2\vec{x}$, so $\lambda^2$ is an eigenvalue of
$\mat{A}^2$ with the \emph{same} eigenvector. \textbf{(b)} If $\mat{A}$ is
invertible, show $\mat{A}\inv\vec{x}=\tfrac{1}{\lambda}\vec{x}$. \textbf{(c)}
Show $(\mat{A}+c\mat{I})\vec{x}=(\lambda+c)\vec{x}$. Conclude that shifting a matrix
by $c\mat{I}$ shifts every eigenvalue by $c$ but leaves the eigenvectors alone.
\end{exercise}

\begin{exercise}[A $3\times 3$ from start to finish]
\label{ch23:ex-3full}
For $\mat{A}=\begin{bmatrix}4 & 0 & 1\\ -2 & 1 & 0\\ -2 & 0 & 1\end{bmatrix}$,
find the characteristic polynomial, all three eigenvalues, and an eigenvector for
each. Check your eigenvalues against the trace and the determinant.
\end{exercise}

\begin{exercise}[Repeated eigenvalue, two outcomes]
\label{ch23:ex-repeat}
Both $\mat{I}=\begin{bmatrix}5 & 0\\ 0 & 5\end{bmatrix}$ and
$\mat{J}=\begin{bmatrix}5 & 1\\ 0 & 5\end{bmatrix}$ have the single eigenvalue
$\lambda=5$ (doubled). For each, find the dimension of the eigenspace for
$\lambda=5$. Explain in words why one of these has a full set of eigenvectors and
the other does not, and connect this to the shear of Section~\ref{ch23:special}.
\end{exercise}

\begin{exercise}[Trace/determinant proof]
\label{ch23:ex-proof}
Prove that for any $2\times 2$ matrix $\mat{A}=\begin{bmatrix}a&b\\c&d\end{bmatrix}$,
the characteristic polynomial equals $\lambda^2-(\tr\mat{A})\lambda+\det\mat{A}$,
and deduce that the two eigenvalues satisfy $\lambda_1+\lambda_2=\tr\mat{A}$ and
$\lambda_1\lambda_2=\det\mat{A}$. Then show, as a consequence, that a real
$2\times 2$ matrix has complex (non-real) eigenvalues exactly when
$(\tr\mat{A})^2<4\det\mat{A}$.
\end{exercise}

\begin{exercise}[The Markov long run -- challenge]
\label{ch23:ex-markov}
Let $\mat{M}=\begin{bmatrix}0.8 & 0.3\\ 0.2 & 0.7\end{bmatrix}$. Its entries are
nonnegative and each column sums to $1$, so it is a \emph{transition matrix}: entry
$M_{ij}$ is the probability of moving from state $j$ to state $i$ in one step, as in
the PageRank box. A process that moves at random between states, with the chance of
each move depending only on the current state, is called a \emph{Markov chain}.
\textbf{(a)} Show $\lambda=1$ is an
eigenvalue and find its eigenvector $\vec{q}$; scale $\vec{q}$ so its entries sum to
$1$. \textbf{(b)} Find the other eigenvalue and confirm $|\lambda|<1$.
\textbf{(c)} Using the eigendecomposition idea of
Section~\ref{ch23:geometry}, argue that for \emph{any} starting probability vector
$\vec{p}_0$, the sequence $\mat{M}^n\vec{p}_0$ converges to $\vec{q}$. Why does the
eigenvalue $1$ ``win,'' and what role does the second eigenvalue play in how fast
you get there?
\end{exercise}

\begin{exercise}[Complex eigenvectors -- challenge]
\label{ch23:ex-complex-evec}
For the $90^\circ$ rotation $\mat{Q}=\begin{bmatrix}0 & -1\\ 1 & 0\end{bmatrix}$
with eigenvalue $\lambda=i$, solve $(\mat{Q}-i\mat{I})\vec{x}=\vec{0}$ over the
complex numbers to find a complex eigenvector. Then take the eigenvalue
$\lambda=-i$ and show its eigenvector is the complex conjugate of the first. (This
conjugate-pairing of complex eigenvectors mirrors the conjugate-pairing of complex
eigenvalues for real matrices.)
\end{exercise}

\begin{exercise}[Build a matrix from eigendata]
\label{ch23:ex-build}
Construct a $2\times 2$ matrix whose eigenvalues are $3$ and $-1$ with
eigenvectors $\colvec{1\\1}$ and $\colvec{1\\-1}$ respectively. \emph{Hint:} you
want $\mat{A}\colvec{1\\1}=3\colvec{1\\1}$ and $\mat{A}\colvec{1\\-1}=-\colvec{1\\-1}$;
set up and solve for the four entries, or anticipate the next chapter and build
$\mat{A}=\mat{S}\mat{\Lambda}\mat{S}\inv$ with the eigenvectors as the columns of
$\mat{S}$.
\end{exercise}
